
%BeginMC
\question
Dinosaurs are not extinct because 

\begin{choices}
\correctchoice birds are part of a monophyletic group with saurischian dinosaurs. %EndChoice
\choice birds are part of a monophyletic group with ornithischian dinosaurs. %EndChoice
\choice the archosaurs is a monophyletic group that contains alligators, pterosaurs, dinosaurs, birds, and (based on recent evidence) turtles.   %EndChoice
\choice the ancestors of alligators form a monophyletic group with other dinosaurs. %EndChoice
\choice the ancestors of alligators branched off at the base of the dinosaur tree. 
\choice None of the above are correct. Dinosaurs actually are extinct.%dontshuffle%EndChoice 
\end{choices}
%EndMC


%BeginMC
\question
The order of appearance of feather types in the fossil record of theropod dinosaurs suggests that the earliest function of feathers was

\begin{choices}
\correctchoice to provide thermal insulation. %EndChoice
\correct for gliding among the tree tops. %EndChoice
\choice to reduce turbulence as the theropods jumped from tree to tree. %EndChoice
\choice to provide balance when the theropods chased prey. %EndChoice
\choice to provide defense against very large predators. %EndChoice
\end{choices}
%EndMC

%BeginMC
\question
Single-celled eukaryotes first appeared in the fossil record about 1.8--2.0 billion years ago. A recent study suggested that single-celled eukaryotes were able to evolve from smaller prokaryotic ancestors because

\begin{choices}
\correctchoice atmospheric oxygen reached concentrations suitable to suit larger organisms. %EndChoice
\choice extinction of many prokaryotes emptied niches that could be filled by descendant eukaryotes. %EndChoice
\choice mass extinction of the Ediacaran fauna removed predation pressure, allowing for larger organism size. %EndChoice
\choice increased numbers of \textit{hox} genes allowed for the greater complexity of single-celled eukaryotes. %EndChoice
\choice Pipit, my cat, willed it to be so.%dontshuffle%EndChoice 
\end{choices}
%EndMC


%BeginMC
\question
In your study of giraffe species, you used genetic comparisons among several morphologically similar populations. Why are genetic tests like this a good way to discover cryptic species?

\begin{choices}
\correctchoice Grouping similar genotypes might reveal populations that do not or cannot interbreed. %EndChoice
\choice Cryptic species will show more genetic differences. %EndChoice
\choice The alleles of a species will be identical among individuals within a population. %EndChoice
\choice They allow you to compare extinct individuals with living individuals. %EndChoice
\choice Cryptic species will show fewer genetic differences. %EndChoice
\end{choices}
%EndMC


%BeginMC
\question
Co-option of a gene occurs when

\begin{choices}
\correctchoice one gene evolves a new function in addition to its original function. %EndChoice
\choice a trait evolves a new function, such as feathers evolving from insulation to sexual selection functions. %EndChoice
\choice two different species mate, creating genetic incompatibility in the hybrids. %EndChoice
\choice individuals from different populations mate, creating a heterozygous advantage. %EndChoice
\choice Two species evolve cooperative behavior (like mutualism), increasing fitness of their offspring. %EndChoice
\end{choices}
%EndMC


%BeginMC
\question
In the state of Hawaii, the most closely related species of \textit{Laupala} crickets occur together on the same island, suggesting they may have evolved by sympatric speciation.  To convince your skeptical science colleagues, you would have to demonstrate that

\begin{choices}
\correctchoice the species have never been distributed allopatrically in the past. %EndChoice
\choice the species are reproductively isolated. %EndChoice
\choice the species occupy different ecological niches. %EndChoice
\choice the species breed at different times. %EndChoice
\choice the species are capable of dispersing to neighboring islands. %EndChoice
\end{choices}
%EndMC


%BeginMC
\question
The type of jaw joint in early mammals that separates them from their reptile-like ancestors is 
\begin{choices}
\correctchoice dentary-squamosal. %EndChoice
\choice quadrate-articular. %EndChoice
\choice dentary-articular. %EndChoice
\choice dentary-quadrate. %EndChoice
\choice articular-squamosal. %EndChoice
\choice quadrate-squamosal. %EndChoice
\end{choices}
%EndMC


%BeginMC
\question
What role(s) does \textit{Bmp4} play in the expression of beak phenotypes?

\begin{choices}
\correctchoice  It is a regulatory gene that determines beak depth and width. %EndChoice
\choice It is a structural protein that contributes to the strength of the beak.%EndChoice
\choice It is a regulatory gene that determines beak length and width. %EndChoice
\choice It is a regulatory gene that inhibits calmodulin, reducing the beak length. %EndChoice
\choice It is a regulatory gene that promotes calmodulin, increasing the beak width. %EndChoice
\choice It is a structural protein that contributes to the curvature of the beak. %EndChoice
\end{choices}
%EndMC


%BeginMC
\question
Which of the following is an example of a postzygotic isolating barrier?

\begin{choices}
\correctchoice The hybrid offspring of two species of crocodiles that can produce normal gametes but cannot obtain a mate due to conflicting behaviors. %EndChoice
\choice The genitalia of a male duck that does not fit properly with the genitalia of females from another population. %EndChoice
\choice Females of one kind of fly that are not attracted to the buzz of another kind of male fly. %EndChoice
\choice Two species of bat that breed in different habitats.%EndChoice
\choice Secretions in the oviduct of a female spider kills sperm from another kind of male after copulation.  %EndChoice
\end{choices}
%EndMC

%BeginMC
\question
Inbreeding results in a higher frequency of \rule{1in}{0.4pt} in a population.  Inbreeding depression occurs because \rule{1in}{0.4pt}.

\begin{choices}
\correctchoice homozygosity; deleterious recessive alleles are expressed more often. %EndChoice
\choice heterozygosity; heterozygotes have lower fitness. %EndChoice
\choice deleterious alleles; individuals with deleterious alleles have high mortality. %EndChoice
\choice heterozygosity; deleterious dominant alleles are expressed more often. %EndChoice
\choice rare alleles; individuals with rare alleles reproduce less often. %EndChoice
\end{choices}
%EndMC
